\chapter{The server}
\label{ch:server}

\cmd{fxmenet-server} is a console program with no window. It hosts plugins
for any number of clients at once, as far as the CPU allows.

\section{Command line}

\begin{lstlisting}
fxmenet-server [--name NAME] [--control-port N] [--udp-port N]
               [--discovery-port N] [--no-discovery]
               [--plugin-path DIR]... [--no-default-paths] [--rescan]
               [--remote-admin]
\end{lstlisting}

\begin{longtable}{@{}p{0.34\textwidth}p{0.6\textwidth}@{}}
\toprule
\textbf{Option} & \textbf{Meaning}\\
\midrule
\endhead
\cmd{--name NAME} & The name shown to clients. Default: the computer's name.\\
\cmd{--control-port N} & TCP port for control. Default 54322.\\
\cmd{--udp-port N} & UDP port for audio. Default 54321.\\
\cmd{--discovery-port N} & UDP port answering \ui{Find}. Default 54320.\\
\cmd{--no-discovery} & Do not answer \ui{Find}. Clients must then type the address.\\
\cmd{--plugin-path DIR} & Also search this folder, for this run only. May be repeated.\\
\cmd{--no-default-paths} & Do not search the standard VST3 folders, for this run only.\\
\cmd{--rescan} & Ignore the plugin cache and examine every plugin again.\\
\cmd{--remote-admin} & Let \cmd{fxmenet-top} change the plugin folders from other machines (section~\ref{sec:admin}).\\
\bottomrule
\end{longtable}

The server logs to its standard output, with a time stamp on each line, so a
service manager (such as systemd's journal) collects it.

\section{Which plugins the server offers}
\label{sec:catalog}

The server offers VST3 \emph{effects}. Instruments are skipped. A plugin
bundle that holds several plugins (a suite) offers each of them.

Each plugin file is examined in a separate process. A plugin that crashes or
hangs while it is loaded is logged, skipped, and does not take the server
down. It stays skipped until the file changes. The results are kept in a
cache file (\file{fxmenet-server-plugins.xml}, in the \file{FX-Mechanics}
folder of the user's configuration directory), keyed by the file's path, size
and date, so only new or changed plugins are examined at the next start. This
examination also notes whether the plugin is made for \fxme\
(chapter~\ref{ch:developers}).

If the same plugin is installed in two folders, the first one found is used.

\section{Plugin folders}
\label{sec:folders}

The server searches three kinds of folders:

\begin{description}
  \item[Default folders.] The standard VST3 folders of the operating system.
    They can be turned off.
  \item[Saved folders.] Folders you add with \cmd{fxmenet-top}
    (chapter~\ref{ch:top}). The server remembers them from one run to the
    next, in \file{fxmenet-server-settings.json} beside the cache file.
  \item[This run's folders.] Those given with \cmd{--plugin-path}. They are
    not remembered.
\end{description}

Folders are searched with their subfolders. After a change of folders, or
when you install a plugin while the server runs, ask for a rescan in
\cmd{fxmenet-top}. The scan runs in the background; the sessions in progress
keep playing. A client sees the new list the next time it connects.

\section{Ports and discovery}

The server listens on three ports (chapter~\ref{ch:network}). \ui{Find}
in the plugin broadcasts a question on the discovery port, and every server
answers with its name, system and number of sessions. Broadcasts normally
stay inside one network segment: across routers or VPNs, type the address.
Only one server can run on a machine with discovery on; a second one says so
in its log and runs without it.

\section{Sessions}

A session starts when a client connects and ends when it says goodbye. If a
client's connection drops (a network hiccup, a DAW that crashes), the server
keeps the session for 10 seconds, so that the client can resume it without
losing its chain.

Each session has its own audio thread. The server asks the operating system
to run it with real-time priority; if the system refuses, the server logs a
warning and runs it as an ordinary thread, which is less reliable under load
(section~\ref{sec:linuxtuning}).

\section{Reading the log}
\label{sec:log}

At startup the log lists the plugins found. Every 10 seconds, for each
session that is working, it prints one line of figures (the numbers
below are an example):

\begin{lstlisting}[language={}]
session 3: 8123 blocks, lost 0, stale 0, from cache 0, queue full 0,
send errors 0, resets 1 | process 412 us, max 10 s 1180 us, max 2310 us |
turnaround max 10 s 1650 us, over 1 block 0, over 2 blocks 0 |
chain latency 256, telemetry frames 240, parameter changes 12
\end{lstlisting}

Counts are since the session started; the values marked \emph{10 s} cover the
last interval.

\begin{longtable}{@{}p{0.27\textwidth}p{0.67\textwidth}@{}}
\toprule
\textbf{Figure} & \textbf{Meaning}\\
\midrule
\endhead
blocks & Blocks processed.\\
lost & Blocks that never arrived in time: the server skipped them (the client conceals them).\\
stale & Late blocks too old to answer, dropped.\\
from cache & Duplicate requests answered from the server's cache of recent replies.\\
queue full & Requests refused because the session could not keep up.\\
send errors & Replies the system refused to send, typically because the send buffer was full (section~\ref{sec:linuxtuning}).\\
resets & Times the client asked to restart the chain's audio (transport start, stop or jump).\\
process & Time spent processing the last block, then the worst in the last 10 seconds, then the worst since the start.\\
turnaround & Time from a request's arrival to its reply leaving, worst of the last 10 seconds.\\
over 1 block, over 2 blocks & Replies that left more than one (two) block durations after the request arrived: the server fell behind.\\
chain latency & Latency of the chain, in samples.\\
\bottomrule
\end{longtable}

A \emph{process} time that gets close to the block's duration (2.9\,ms for 128
samples at 44.1\,kHz) means the server is running out of CPU for that
session.

\section{Running the server as a service}

Below is an example of a \cmd{systemd} unit for Linux. Running the server as a service is a possibility that has not been tested yet, so treat the example as a starting point. Feedback is welcome, and so are equivalents for other systems (launchd on macOS, a Windows service).

\begin{lstlisting}[language={}]
[Unit]
Description=FxmeNet server
After=network-online.target
Wants=network-online.target

[Service]
ExecStart=/usr/local/bin/fxmenet-server
User=fxmenet
Group=audio
Restart=on-failure
LimitRTPRIO=95
LimitMEMLOCK=infinity

[Install]
WantedBy=multi-user.target
\end{lstlisting}

\section{Letting a remote \cmd{fxmenet-top} change the folders}
\label{sec:admin}

\cmd{fxmenet-top} can watch a server from any machine of the network. It may
change the plugin folders only when it connects from the server's own
machine, unless you start the server with \cmd{--remote-admin}. Because the
network is not authenticated, use this option only on a network you
trust.
